% ============================================
% 西南大学本科毕业论文（设计）
% 编译：双击 compile.bat
% ============================================
\documentclass[12pt,a4paper]{swuthesis}

\usepackage{graphicx}
\usepackage{mhchem}
\usepackage{siunitx}

\addbibresource{reference.bib}

% ============================================
% 以下信息按实际情况修改
% ============================================
\title{论文题目（二号宋体，不超过20字）}
\swucollege{学院名称}
\swumajor{专业名称}
\swuauthorid{学号}
\swuauthor{姓名}
\swuadvisor{指导教师}
\swudate{2029年5月}
\swuentitle{English Title}
\swuenname{Name}
\swuencollege{College Name}
\swuabstract{%
这里是中文摘要，约300字，简要说明研究目的、方法、结果和结论。}
\swuenabstract{%
English abstract here, about 250 words.}
\swukeywords{关键词A；关键词B；关键词C；关键词D}
\swuenkeywords{KeywordA; KeywordB; KeywordC; KeywordD}

% ============================================
\begin{document}

\maketitle           % 封面
\makedeclaration     % 独创声明
\makecontents        % 目录
\makeabstract        % 中英文摘要

% ==================== 第1章 导论 ====================
\section{导论}

\subsection{研究背景与意义}

正文宋体小四号，1.5倍行距。首行自动缩进。此处撰写研究背景、意义、目的、现状等。

引用示例：PDMS具有优异的疏水性能\citeup{zhang2023}。Wenzel模型描述了粗糙表面的润湿行为\citeup{wenzel1936}。

\subsection{国内外研究现状}

\subsubsection{国外研究现状}

国外研究……

\subsubsection{国内研究现状}

国内研究……

\subsection{研究内容}

本文主要研究……

% ==================== 第2章 实验部分 ====================
\newpage
\section{实验部分}

\subsection{试剂与仪器}

\subsubsection{实验试剂}

实验试剂如\ref{tab:reagents}所示。

\begin{table}[H]\centering
    \caption{\zihao{5} 实验试剂}
    \addtocounter{table}{-1}
    \renewcommand{\tablename}{Table}
    \caption{\zihao{5} Experimental reagents}
    \renewcommand{\tablename}{表}
    \label{tab:reagents}
    \begin{tabular}{cccc}\toprule
        试剂 & 规格 & 厂家 & 用途 \\\midrule
        PDMS & AR & 道康宁 & 基底 \\
        固化剂 & AR & 道康宁 & 交联 \\\bottomrule
    \end{tabular}
\end{table}

\subsubsection{实验方法}

PDMS薄膜制备：……

接触角测量使用Young方程：
\begin{equation}
    \cos\theta = \frac{\gamma_{SV} - \gamma_{SL}}{\gamma_{LV}}
    \label{eq:young}
\end{equation}

% ==================== 第3章 结果与讨论 ====================
\newpage
\section{结果与讨论}

\subsection{润湿性能分析}
结果如\ref{fig:ca}所示。

\begin{figure}[H]\centering
    \fbox{\parbox{0.5\textwidth}{\centering\vspace{3cm}插入接触角图片}}
    \caption{\zihao{5} 水接触角测量结果}
    \addtocounter{figure}{-1}
    \renewcommand{\figurename}{Fig.}
    \caption{\zihao{5} Water contact angle results}
    \renewcommand{\figurename}{图}
    \label{fig:ca}
\end{figure}

\subsection{讨论}
分析表明……

% ==================== 第4章 结论 ====================
\newpage
\section{结论}

\begin{enumerate}
    \item 结论一。
    \item 结论二。
    \item 结论三。
\end{enumerate}

% ==================== 参考文献 ====================
\makereference

% ==================== 致谢 ====================
\makethanks{感谢导师悉心指导。感谢课题组同窗帮助。感谢家人支持。}

% ==================== 附录（需要时取消注释） ====================
% \makeappendix
% \subsection{附录A  实验数据}
% 附录内容……

\end{document}
